\section{AI Safety: descomponer un problema enorme en un problema de ingeniería ejecutable}

\subsection{El marco por capas de Jensen}

En la entrevista, Jensen descompone AI safety en niveles ejecutables, un enfoque coherente con la ingeniería de seguridad de la aviación, el automóvil y el control industrial:
\begin{itemize}
    \item Riesgo a nivel de modelo: razonamiento incorrecto, alucinaciones, entradas adversarias.
    \item Riesgo a nivel de sistema: fallas de hardware, interrupciones de red, fallas del enlace de control.
    \item Riesgo a nivel de uso: fraude, deepfake, suplantación de identidad, manipulación de información.
\end{itemize}

\begin{importantbox}{Una visión de la seguridad orientada a la ingeniería}
Él enfatiza que muchos riesgos no son «la IA quiere hacer el mal», sino «el sistema no fue diseñado ni validado correctamente». Por eso, las medidas de gobernanza deben priorizarse en:
\begin{itemize}
    \item arquitectura redundante (Redundancy)
    \item protección contra fallas (Fail-safe)
    \item pruebas continuas y monitoreo en línea
\end{itemize}
\end{importantbox}

\subsection{La seguridad no es una «métrica del modelo», sino una «cadena de responsabilidad de extremo a extremo»}

Para la conducción autónoma, la asistencia médica y la robótica industrial, la exactitud del modelo es solo una métrica de entrada. La seguridad realmente utilizable depende de la cadena de responsabilidad completa: recopilación de datos, entrenamiento del modelo, control del despliegue, toma de control humana y análisis posterior de incidentes.

\begin{warningbox}{Error común: reducir AI safety a «si se prohíbe o no cierto modelo»}
Si la gobernanza de la seguridad se centra solo en la prohibición a nivel de modelo, pasará por alto una gran cantidad de riesgos a nivel de sistema. La seguridad de los sistemas realmente de alto riesgo proviene de la disciplina de ingeniería y la trazabilidad de los procesos, no solo de la elección del modelo.
\end{warningbox}

\subsection{Resumen de la sección}
La visión sobre safety en la entrevista puede resumirse como «convertir un problema complejo en un problema de ingeniería verificable». Esta metodología es más adecuada para la etapa de implementación industrial.

